\chapter{NOTAȚII MATEMATICE. REGULI ORTOGRAFICE}

\section{Mărimi și unități de măsură}

În cadrul aceleiași lucrări, nu se admite ca o mărime să fie notată cu mai multe simboluri. De asemenea, se va evita pe cât posibil ca mărimi diferite să fie notate cu același simbol.

Unitățile de măsură folosite trebuie să aparțină Sistemului Internațional de Unități de Măsură (SI) și, tot conform acestuia, se vor adopta simbolurile lor.


\section{Notații folosite în formulele matematice}

\subsection{Reguli generale}

Simbolul pentru scădere (minus) trebuie să fie linia lungă (\textbf{--}, \textit{En dash}), nu cratima (-).

Valorile numerice (mărimi, indici, exponenți etc.) se vor scrie drept (nu \textbf{bold}/\textit{italic}).

Mărimile se vor nota cu simboluri scrise \textit{italic} (indiferent dacă se folosesc litere ale alfabetului latin sau litere grecești), iar valoarea numerică a mărimii și unitatea de măsură se vor scrie drept, cu spațiu între ele – ex.: $f_e$ = 16 kHz, $\sigma^2$ = 1.

În cazul utilizării unei unități de măsură separat de o valoare numerică (ex.: etichetarea axelor figurilor), unitatea de măsură se va scrie între paranteze drepte; ex.: $f$ [Hz], $|S(e^{j\omega})|$ [dB].

Mărimile de natură vectorială se vor scrie \textbf{bold} și drept; vectorii vor fi scriși cu minuscule, matricele cu majuscule – ex.: \textbf{x}\textsubscript{\textit{n}} = [\textit{F}\textsubscript{0}; \textit{F}\textsubscript{1}; ...; \textit{F}\textsubscript{\textit{D}-1}], \textbf{X} = [\textbf{x}\textsubscript{0}, \textbf{x}\textsubscript{1}, ..., \textbf{x}\textsubscript{\textit{N}-1}].

Ecuațiile vor fi centrate, pe rând nou. Între ecuații și text se va lăsa o distanță de 1 rând liber (12 pt) deasupra și dedesubtul lor. Explicațiile notațiilor dintr-o ecuație se recomandă să fie scrise în paragraful în care este referită ecuația.

Ecuațiile vor fi numerotate în continuare, pe capitole (ex.: (\ref{Eq_freq}) ). Numărul curent al ecuației se va scrie între paranteze rotunde, plasat în dreapta, la marginea formatului textului.

Toate ecuațiile se vor numerota și trebuie referite în text înainte de scrierea acestora.

\subsection{Alte reguli}

	Pentru formatul fonturilor:
 \begin{enumerate}
     \item se vor mai scrie \textit{italic}:
     \begin{enumerate}
         \item variabile \textbf{--} ex.: \textit{x}, \textit{y}, \textit{z};
         \item funcții oarecare \textbf{--} ex.: \textit{f}, \textit{g}, \textit{f}(x), \textit{v}(n);
         \item indici/exponenți care reprezintă la rândul lor simbolul unei mărimi \textbf{--} ex.: $e^{j\omega}$;
         \item cuvinte în limbi străine care \textbf{nu} sunt incluse în DEX \textbf{--} ex.: funcție \textit{kernel}; pe cât posibil, însă, \textbf{se va încerca traducerea corespunzătoare a acestora}.\\
     \end{enumerate} 
     \item se vor mai scrie drept:
     \begin{enumerate}
         \item funcțiile matematice consacrate \textbf{--} ex.: funcții trigonometrice (sin, cos, tanh etc.), funcții logaritm (log, lg, ln);
	\item notația d pentru derivare  \( \displaystyle \left( \frac{\textnormal{d}v}{\textnormal{d}t} \right) \) și simbolul pentru derivare parțială ($\partial$);
	\item simbolurile pentru limită (lim), constant (ct.), integrală ($\int$), unitatea imaginară (i sau j), constante (e, $\pi$ etc.);
	\item indici (inferiori sau superiori) care provin din prescurtări – ex.: min, max, inf, sup etc.\\
     \end{enumerate}
 \end{enumerate}

În finalul acestui capitol este ilustrat un exemplu prin (\ref{Eq_freq}), aceasta fiind condiția lui Nyquist asociată teoremei eșantionării, unde $f_e$ este frecvența de eșantionare, iar $f_m$ reprezintă frecvența maximă din componența semnalului.

\begin{equation}
    f_e > 2 \cdot f_m
    \label{Eq_freq}
\end{equation}

\section{Reguli ortografice}

Pauza \textbf{nu se folosește înainte, ci după} punct (.), virgulă (,), două puncte (:), punct şi virgulă(;), semnul exclamării (!), semnul întrebării (?).
 
În cazul citatelor unde se folosesc ghilimele, pauza se folosește \textbf{înaintea} ghilimelelor de la începutul citatului (nu între ghilimele și primul cuvânt al citatului), iar în cazul ghilimelelor de la sfârșitul citatului pauza se folosește \textbf{după} ghilimele. Aceeași regulă este valabilă în cazul parantezelor.\\

Pentru cuvintele în limbi străine care \textbf{nu} sunt incluse în DEX, sufixul formei de plural și/sau articolul hotărât se separă de cuvânt prin cratimă \textbf{--} ex.: \textit{kernel}-uri, \textit{transformer}-ele.

Pentru cuvintele în limbi străine care au fost incluse în DEX (fără a li se adapta forma), sufixul formei de plural și/sau articolul hotărât \textbf{nu} se separă prin cratimă dacă forma de singular a cuvântului se termină cu o \textbf{consoană} \textbf{--} ex.: patternuri, kurtosisul.



